\begin{otherlanguage}{spanish}
\chapter*{Resumen}
\addcontentsline{toc}{chapter}{Resumen}

La evidencia de la materia oscura proviene de numerosas pruebas independientes a distintas escalas. Las curvas de rotación galácticas, los movimientos de las galaxias dentro de los cúmulos, las lentes gravitatorias y las anisotropías del fondo cósmico de microondas requieren todas ellas una forma de materia que supera a los bariones ordinarios aproximadamente en un factor cinco y que interactúa con la luz solo a través de la gravedad. Sin embargo, todas estas pruebas son de naturaleza gravitatoria y dicen poco sobre qué es realmente la materia oscura. Si está formada por partículas masivas de interacción débil, podría aniquilarse o desintegrarse en estados del Modelo Estándar, y los rayos gamma son mensajeros especialmente atractivos: al no ser desviados por los campos magnéticos, llevan consigo la firma espectral de la interacción que los produjo. Durante más de quince años, el Large Area Telescope a bordo del telescopio espacial de rayos gamma Fermi (Fermi-LAT) ha cartografiado el cielo del GeV, donde se esperan estas firmas, y sus medidas constituyen el núcleo de esta tesis.

En esta tesis defendemos que el progreso en las búsquedas indirectas de materia oscura depende hoy de métodos estadísticos y de aprendizaje automático capaces de recuperar señales en entornos dominados por el ruido. El trabajo sigue un arco que va del Centro Galáctico hacia el exterior, hasta la red cósmica, siguiendo la señal esperada en regímenes de relación señal-ruido progresivamente menor. Para ello recurrimos a métodos estadísticos tradicionales, al modelado generativo, a los métodos de verosimilitud implícita y al aprendizaje profundo, aplicados a los datos de Fermi-LAT y a predicciones para instrumentos futuros.

El Centro Galáctico alberga el Exceso del Centro Galáctico (GCE, por sus siglas en inglés), un excedente de fotones GeV que podría señalar la aniquilación de materia oscura o una población no resuelta de púlsares de milisegundo (MSPs). Para poner a prueba la hipótesis de los púlsares, medimos la función de luminosidad en rayos gamma de los MSPs en los cúmulos globulares de la Vía Láctea, y encontramos una luminosidad media $\langle L_\gamma\rangle \sim (1\text{--}8)\times10^{33}$ erg/s entre 0.1 y 100 GeV y una dispersión log-normal ancha, $\sigma_L \sim 1.4\text{--}2.8$. Si el GCE hubiera sido producido por púlsares con estas propiedades, Fermi-LAT ya habría resuelto individualmente entre 17 y 37 de ellos, mientras que solo se conocen tres candidatos --- una tensión que reaviva el debate sobre si el exceso procede de púlsares o de materia oscura.

Nos dirigimos después a objetivos más débiles: subhalos de materia oscura potencialmente ocultos entre las fuentes de Fermi-LAT que carecen de contrapartida astrofísica. Construimos el que es, hasta donde sabemos, el primer modelo generativo de estas fuentes no asociadas a latitudes superiores a diez grados: una mezcla estadística de clases astrofísicas conocidas y una posible componente de subhalos, entrenada directamente sobre esta población y corregida por posibles desplazamientos de covariables y de probabilidades a priori.
No encontramos ninguna contribución significativa de subhalos y establecimos límites superiores al 95\% de confianza sobre la sección eficaz de aniquilación en el canal $b\bar{b}$ para masas de materia oscura entre 10 GeV y 1 TeV.

Por debajo del umbral de detección de Fermi-LAT, las fuentes individuales se emborronan entre sí, y el objeto de estudio pasa a ser la población de fuentes de rayos gamma en su conjunto --- tanto las astrofísicas como las posibles fuentes de materia oscura. Entrenamos una red neuronal convolucional con alrededor de $10^6$ mapas sintéticos del cielo de Fermi-LAT para reconstruir la distribución de recuento de fuentes $dN/dS$ de las fuentes extragalácticas a partir de catorce años de datos entre 1 y 10 GeV. La reconstrucción alcanza flujos un factor 50 por debajo del umbral de Fermi-LAT, concuerda con los catálogos existentes en el régimen resuelto y se extiende como $dN/dS \propto S^{-2}$ hasta aproximadamente $5\times10^{-12}$ cm$^{-2}$ s$^{-1}$.
Sobre esta distribución desarrollamos un marco de catalogación probabilística que compara el cielo observado con simulaciones y asigna a cada dirección una probabilidad de albergar una fuente sub-umbral. Este marco identifica aproximadamente un 50\% más de direcciones candidatas a fuente que el catálogo 4FGL-DR3, y proporciona un catálogo público particularmente adecuado para estudios de correlación cruzada.

A las escalas más grandes, la señal de aniquilación sigue la propia distribución de materia, de modo que correlacionar las anisotropías de rayos gamma con las posiciones de las galaxias puede separar una componente cosmológica de materia oscura de los fondos isótropos. Predecimos la sensibilidad del Cherenkov Telescope Array Observatory (CTAO) a esta señal. Con un catálogo denso de galaxias a bajo corrimiento al rojo como 2MASS y unas 50 horas de observación, la correlación cruzada alcanza sensibilidades a la materia oscura tanto en aniquilación como en desintegración comparables a las de los análisis de galaxias enanas y cúmulos de galaxias.

En conjunto, estos resultados trazan un camino coherente a través del cielo de rayos gamma, desde la señal anticipada más brillante hasta la débil huella de la red cósmica. Ninguna de estas búsquedas devuelve una detección confirmada, y la identidad de la partícula de materia oscura sigue siendo un problema abierto. Los próximos pasos de la detección indirecta, defendemos, requerirán mejores instrumentos, más datos y herramientas estadísticas más afiladas para leerlos, avanzando todos a la vez. Esta tesis contribuye al tercero de estos frentes, trabajando en nuevas herramientas estadísticas para el campo: modelos de poblaciones que hacen estadísticamente accesibles las fuentes sub-umbral y marcos de inferencia robustos frente al desajuste entre simulaciones y datos reales, junto con catálogos probabilísticos y análisis de correlaciones cruzadas. La esperanza es que este trabajo de base ayude a extender estas búsquedas a medida que los sondeos actuales de rayos gamma sigan creciendo y los observatorios de próxima generación entren en funcionamiento.
\end{otherlanguage}
